\uuid{v9PP}
\exo7id{6057}
\titre{exo7 6057}
\auteur{queffelec}
\organisation{exo7}
\datecreate{2011-10-16}
\isIndication{false}
\isCorrection{false}
\chapitre{Espace vectoriel normé}
\sousChapitre{Espace vectoriel normé}
\cpge{mpsi-mp2i.espaces-prehilbertiens.b, mp-mpi.topologie-espaces-normes.a, pc.espaces-vectoriels-normes.a, psi.espaces-vectoriels-normes.a, mp-mpi.topologie-espaces-normes.c, pc.espaces-vectoriels-normes.c, psi.espaces-vectoriels-normes.c}
\module{Topologie}
\niveau{L3}
\difficulte{}

\contenu{
\texte{
On considère dans ${\Rr}^2$, les deux applications
$$n((x,y))=\sup_{t\in [0,1]}\vert x+ty\vert,$$
$$m((x,y))=\int_0^1\vert x+ty\vert\ dt.$$
}
\begin{enumerate}
    \item \question{Montrer que $n$ et $m$ définissent deux normes sur ${\Rr}^2$.}
    \item \question{Dessiner les boules unités fermées associées, et trouver des constantes
effec\-tives $A$,$B$, telles que $A\ n((x,y))\leq m((x,y))\leq B\ n((x,y))$ pour
tout $(x,y)\in {\Rr}^2$.}
\end{enumerate}
}
